\section{Scope and notation}
The bounded study \cite{Finite} applies exact categorification criteria
to the rank-four fusion-ring census through multiplicity 128 and gives
unitary models for the fifteen known categorifiable rings.
The present paper asks which parts of that classification can be made
independent of the multiplicity. The finite theorem is an input here,
not a consequence of the unrestricted reductions: there is no circular
use of the two papers.

We use the same basis convention $1,b_1,b_2,b_3$,
$(N_i)_{jk}=N_{ij}^k$, Frobenius--Perron dimensions $d_i$, and
$D=\sum_i d_i^2$. A proper nontrivial fusion subring has rank two or
three. A character field is generated by the values of one character;
it need not be the field generated by all formal codegrees. We distinguish
these fields throughout. For the complete realization table and the fixed
labels $F_{m,j}$ we refer to \cite{Finite}.

The following are used from that paper: the standard arithmetic inputs,
Galois reduction of dimension characters, the low-induction Gram identity,
the two-moment bound, and positivity of squared central dimensions.
Their uses below are explicitly cited. The dimension and parameter bounds
proved here make each subsequent finite computation exhaustive. The two
supporting scripts for the unrestricted Fibonacci and quadratic-character
arguments are supplied separately from the bounded classification verifier.

The conclusions concern existence of a categorification of a given based
ring. Equality of pivotal, spherical and pseudo-unitary \emph{existence}
does not assert that every category is pivotal or every spherical
realization is pseudo-unitary without Galois conjugation. Nor do the results
below prove an unrestricted general or pivotal classification of every
rank-four fusion ring.
